\chapter{関連研究}

\section{sEMGを用いたジェスチャ分類}
これまでに,表面筋電位（sEMG）を用いた手指・手首ジェスチャ分類に関する研究が多数報告されている.一般的な処理フローは,信号取得,セグメンテーション,特徴量抽出,および分類から構成される.特徴量としては,時間領域特徴量,周波数領域特徴量,およびそれらを組み合わせた手法が用いられてきた.

時間領域特徴量は,平均絶対値（MAV）,波形長（WL）,ゼロ交差数（ZC）などに代表され,計算量が少なくリアルタイム処理に適している.一方,周波数領域やウェーブレット領域の特徴量は高い表現力を有するものの,計算コストが高く,組込み環境での利用には制約がある.そのため,近年のリアルタイム応用やエッジデバイスを対象とした研究では,時間領域特徴量を中心とした設計が主流となっている.

分類器としては,線形判別分析（LDA）,サポートベクタマシン（SVM）,人工ニューラルネットワーク（ANN）などが用いられてきた.特にANNは,非線形なsEMG信号の特徴を表現できる点から,多くの研究で採用されている.

\section{単一チャネルsEMGの研究}
単一チャネルsEMG信号を用いたジェスチャ分類は,ハードウェア構成の簡素化や低コスト化の観点から注目されているが,情報量の制約により分類性能の確保が課題とされてきた.

Aroziら\cite{arozi2020}は,単一チャネルsEMGから16次元の時間領域特徴量を抽出し,主成分分析（PCA）によって次元削減を行った後,人工ニューラルネットワーク（ANN）を用いて9種類の手指動作を分類する手法を提案している.この研究では,多チャネルを用いずとも一定の分類精度が得られることを示した一方,PCAによる前処理が必要であり,リアルタイム組込み実装の観点では処理負荷が増加する可能性がある.

また,Ibrahimら\cite{ibrahim}は,単一チャネルsEMGを用いたジェスチャ分類器をESP32マイクロコントローラ上に直接実装する研究を報告している.時間領域特徴量10次元を入力とし,1隠れ層のANNによる分類を行うことで,低遅延かつ高精度なリアルタイム推論を実現している.この研究は,組込みデバイス単体での推論実装の実現性を示した点で重要であるが,学習自体はオフライン環境で行われており,モデルの更新や個人適応については十分に検討されていない.

\section{特徴量設計およびニューラルネットワークモデル}
sEMG分類における特徴量設計は,モデル性能および実装可能性に大きく影響する要素である.近年では,生波形を直接入力とする畳み込みニューラルネットワーク（CNN）を用いた手法も提案されているが,この場合,高い計算資源とメモリ容量が必要となる.

一方,自己回帰（AR）モデル係数を含む時間領域特徴量は,信号の時間構造を要約した低次元表現であり,軽量な分類器と組み合わせることで高い実用性を示すことが報告されている.AR係数は時系列信号としての局所構造を統計的に表現するものであり,畳み込み処理を前提としたCNNとの親和性は高くない.このため,AR係数を含む特徴量設計では,全結合型ニューラルネットワーク（Dense層）による分類が適していると考えられる.

ANNは,入力層,隠れ層,出力層から構成されるフィードフォワード型モデルであり,少数の隠れ層でも非線形分類問題を表現可能である.特に1隠れ層のANNは,計算量と表現力のバランスに優れ,組込み・オンデバイス推論において広く採用されている.

\section{本研究の位置づけ}
本研究では,先行研究で提案されている複数の特徴量設計およびモデル構成について,自身の研究で用いるハードウェアおよびソフトウェア環境において再現実験を行った.その結果,時間領域特徴量10次元（AR係数4次元を含む）をそのまま入力とし,1層の全結合層を介して4クラスを出力するANN構成が,分類精度および処理速度の両面で最も優れていることを確認した.

以上の検討結果を踏まえ,本研究ではこの軽量なANNモデルを採用するとともに,モデルをアップデート可能な構成とし,オンデバイス学習による分類精度向上の可能性を検証することを目的とする.すなわち,既存研究で主に扱われてきた「推論のオンデバイス化」に加えて,「学習を含めたオンデバイス適応」に着目する点に本研究の特徴がある.
